\ea\label{ex:text2-1}
O ATOVO\\
\gll o	atovo\\
     \textsc{art}3.\textsc{sg}	sago.palm\\
\glt `THE SAGO PALM'
\z

\ea\label{ex:text2-2}
O naono vai o atovo na nebaneba komana bata ni riori.\\
\gll o	naono	vai	o	atovo	na	neba{$\sim$}neba	komana	bata	ni	ri=ori\\
     \textsc{art}3.\textsc{sg}	palm	\textsc{dem}6	\textsc{art}3.\textsc{sg}	sago.palm	\textsc{tm}2	\textsc{red}{$\sim$}plant	self	\textsc{cont}	\textsc{appl}	3\textsc{pl}.\textsc{ipfv}=3\textsc{pl}.\textsc{pron}\\
\glt `The sago palm is planted.'\footnote{lit. `This tree, the sago palm, they plant it themselves.'}
\z

\ea\label{ex:text2-3}
Ei kou o naono bona, eve to vaa'ave riori bona maa inu teori.\\
\gll ei~kou	o	naono	bona	eve	to	vaa'ave	ri=ori	bona	maa	inu	te=ori\\
     because	\textsc{art}3.\textsc{sg}	tree	\textsc{ana}.\textsc{sg}	3\textsc{sg}.\textsc{pron}	\textsc{rel}	thatch	3\textsc{pl}.\textsc{ipfv}=3\textsc{pl}.\textsc{pron}	\textsc{obj}.\textsc{art}2.\textsc{sg}	\textsc{plm}	house	\textsc{prep}=3\textsc{pl}.\textsc{pron}\\
\glt `Because it is this palm that they use for thatching their houses.'\footnote{lit. `Because this tree, it (is) what they thatch their houses with.'}
\z

\ea\label{ex:text2-4}
O atovo mee, be vua, a vua nae repaa kuu, ore paa pura komana vahaa.\\
\gll o	atovo	mee	be	vua	a	vua	nae	re=paa	kuu	o=re	paa	pura	komana	vahaa\\
     \textsc{art}3.\textsc{sg}	sago.palm	also	when	bear.fruit	\textsc{art}2.\textsc{sg}	fruit	3\textsc{sg}.\textsc{poss}.\textsc{pron}	\textsc{consec}=\textsc{tm}3	fall	3.\textsc{pron}=\textsc{consec}	\textsc{tm}3	grow	self	again\\
\glt `Also,  when the sago palm bears nuts, (and) a nut falls down, then it grows by itself again.'\footnote{lit. `Also, the sago palm, when it bears fruit, ...'}
\z

\ea\label{ex:text2-5}
Beve sun rova nana tea hatana na pura vakis nana.\\
\gll beve	sun	rova	nana	te=a	hatana	na	pura	vakis	nana\\
     when.3\textsc{sg}.\textsc{pron}	stand	even	3\textsc{sg}.\textsc{ipfv}	\textsc{prep}=\textsc{art}2.\textsc{sg}	swamp	\textsc{tm}2	grow	still	3\textsc{sg}.\textsc{ipfv}\\
\glt `Even when it is standing in the swamp, it still grows.'
\z

\ea\label{ex:text2-6}
KOSI ATOVO\\
\gll kosi	atovo\\
     cut.off.leaves	sago.palm\\
\glt `CUTTING OFF THE LEAVES OF THE SAGO PALM'
\z

\ea\label{ex:text2-7}
Bean tau kosi o atovo, taba momohu toro paku an, ean toro tagava roho o atovo.\\
\gll be=an	tau	kosi	o	atovo	taba	momohu	toro	paku	an	e-an	toro	tagava	roho	o	atovo\\
     when=2\textsc{sg}.\textsc{pron}	\textsc{tm}5	cut.off.leaves	\textsc{art}3.\textsc{sg}	sago.palm	thing	first	\textsc{tm}7	do	2\textsc{sg}.\textsc{pron}	\textsc{h}-2\textsc{sg}.\textsc{pron}	\textsc{tm}7	fell	before	\textsc{art}3.\textsc{sg}	sago.palm\\
\glt `When you are about to cut off the leaves of a sago palm, the first thing you must do, you must fell a sago palm.'
\z

\ea\label{ex:text2-8}
Ean repaa mosi hapehape o atovo.\\
\gll e-an	re=paa	mosi	hape{$\sim$}hape	o	atovo\\
     \textsc{h}-2\textsc{sg}.\textsc{pron}	\textsc{consec}=\textsc{tm}3	cut	\textsc{red}{$\sim$}cut.off.branch	\textsc{art}3.\textsc{sg}	sago.palm\\
\glt `Then you cut off the branches of the sago palm.'
\z

\ea\label{ex:text2-9}
Ean repaa kosi amaa paka atovo kahi bono sinoo vaa tobina.\\
\gll e-an	re=paa	kosi	a=maa	paka	atovo	kahi	bono	sinoo	vaa	tobina\\
     \textsc{h}-2\textsc{sg}.\textsc{pron}	\textsc{consec}=\textsc{tm}3	cut.off.leaves	\textsc{art}2.\textsc{sg}=\textsc{plm}	leaf	sago.palm	from	\textsc{obj}.\textsc{art}3.\textsc{sg}	bone	\textsc{lnk}	middle\\
\glt `Then you cut the sago palm leaves off the stalk in the middle.'\footnote{lit. `the bone in the middle'}
\z

\ea\label{ex:text2-10}
Murinae ean repaa varomo nia maa paka atovo.\\
\gll muri-nae	e-an	re=paa	varomo	ni=a	maa	paka	atovo\\
     after-3\textsc{sg}.\textsc{poss}.\textsc{pron}	\textsc{h}-2\textsc{sg}.\textsc{pron}	\textsc{consec}=\textsc{tm}3	stack.up	\textsc{appl}=\textsc{art}2.\textsc{sg}	\textsc{plm}	leaf	sago.palm\\
\glt `After that, you stack up the sago palm leaves.'
\z

\ea\label{ex:text2-11}
Bean varomo vakavara, ean repaa kisi vatataunun riori.\\
\gll be=an	varomo	va-kavara	e-an	re=paa	kisi	va-ta{$\sim$}ta-unun	ri=ori\\
     when=2\textsc{sg}.\textsc{pron}	stack.up	\textsc{advr}-complete	\textsc{h}-2\textsc{sg}.\textsc{pron}	\textsc{consec}=\textsc{tm}3	tie	\textsc{mult}-\textsc{red}{$\sim$}\textsc{acaus}-gather	3\textsc{pl}.\textsc{objm}=3\textsc{pl}.\textsc{pron}\\
\glt `When you have stacked them up, you tie them together.'\footnote{lit. `When you have completely stacked up, ...'}
\z

\ea\label{ex:text2-12}
Ean repaa taneo mana tea dee voosu amaa kisi atovo vo vaan.\\
\gll e-an	re=paa	taneo	mana	tea	dee	voosu	a=maa	kisi	atovo	vo	vaan\\
     \textsc{h}-2\textsc{sg}.\textsc{pron}	\textsc{consec}=\textsc{tm}3	start	eventually	\textsc{compl}1	carry	go.home	\textsc{art}2.\textsc{sg}=\textsc{plm}	bundle	sago.palm	\textsc{goal}	village\\
\glt `Then you evenually start carrying the bundles of sago palm (leaves) home to the village.'
\z

\ea\label{ex:text2-13}
Bea kavara na maa kisi atovo tei vaan,\\
\gll be=a	kavara	n=a	maa	kisi	atovo	tei	vaan\\
     when=\textsc{art}2.\textsc{sg}	all	3\textsc{sg}.\textsc{poss}=\textsc{art}2.\textsc{sg}	\textsc{plm}	bundle	sago.palm	be.there	village\\
\glt `When all the bundles of sago palm (leaves) are in the village,'
\z

\ea\label{ex:text2-14}
ean repaa taneo tea siri sinoo\\
\gll e-an	re=paa	taneo	tea	siri	sinoo\\
     \textsc{h}-2\textsc{sg}.\textsc{pron}	\textsc{consec}=\textsc{tm}3	start	\textsc{compl}1	tear.off	bone\\
\glt `you start tearing off the bones,'
\z

\ea\label{ex:text2-15}
o sinoo vaa tobina no paka atovo to siri noman.\\
\gll o	sinoo	vaa	tobina	n=o	paka	atovo	to	siri	nom=an\\
     \textsc{art}3.\textsc{sg}	bone	\textsc{lnk}	middle	3\textsc{sg}.\textsc{poss}=\textsc{art}3.\textsc{sg}	leaf	sago.palm	\textsc{rel}	tear.off	\textsc{ipfv}=2\textsc{sg}.\textsc{pron}\\
\glt `it is the bone in the middle of the sago palm leaf that you tear off (i.e. the midrib).'
\z

\ea\label{ex:text2-16}
Bean siri vakavara a peha kisi atovo,\\
\gll be=an	siri	va-kavara	a	peha	kisi	atovo\\
     when=2\textsc{sg}.\textsc{pron}	tear.off	\textsc{advr}-complete	\textsc{art}2.\textsc{sg}	one2.\textsc{sg}	bundle	sago.palm\\
\glt `When you have finished tearing off (the midribs of) one bundle of sago palm leaves,'
\z

\ea\label{ex:text2-17}
ean repaa vaoromo vaha bana amaa paka atovo.\\
\gll e-an	re=paa	vaoromo	vahabana	a=maa	paka	atovo\\
     \textsc{h}-2\textsc{sg}.\textsc{pron}	\textsc{consec}=\textsc{tm}3	stack.up	again	\textsc{art}2.\textsc{sg}=\textsc{plm}	leaf	sago.palm\\
\glt `you stack up the sago palm leaves again.'
\z

\ea\label{ex:text2-18}
Murinae ean repaa kaku eve, a vaamatatopo vaa tea huha.\\
\gll muri-nae	e-an	re=paa	kaku	eve	a	vaa-matatopo	vaa	tea	huha\\
     after-3\textsc{sg}.\textsc{poss}.\textsc{pron}	\textsc{h}-2\textsc{sg}.\textsc{pron}	\textsc{consec}=\textsc{tm}3	fold	3\textsc{sg}.\textsc{pron}	\textsc{art}2.\textsc{sg}	\textsc{caus}-be.prepared	\textsc{lnk}	\textsc{compl}1	sew\\
\glt `After that you fold them in preparation for sewing (the thatch).'
\z

\ea\label{ex:text2-19}
O SIKIRI\\
\gll o	sikiri\\
     \textsc{art}3.\textsc{sg}	midrib\\
\glt `THE MIDRIB'
\z

\ea\label{ex:text2-20}
Amaa sikiri to paa siri vuru an tea maa paka atovo,\\
\gll a=maa	sikiri	to	paa	siri	vuru	an	te=a	maa	paka	atovo\\
     \textsc{art}2.\textsc{sg}=\textsc{plm}	midrib	\textsc{rel}	\textsc{tm}3	tear	earlier	2\textsc{sg}.\textsc{pron}	\textsc{prep}=\textsc{art}2.\textsc{sg}	\textsc{plm}	leaf	sago.palm\\
\glt `The midribs that you have earlier torn off the sago palm leaves'
\z

\ea\label{ex:text2-21}
eve to pasi huha atovo vahaa ni noman.\\
\gll eve	to	pasi	huha	atovo	vahaa	ni	nom=an\\
     3\textsc{sg}.\textsc{pron}	\textsc{rel}	\textsc{tm}4	sew	sago.palm	again	\textsc{appl}	\textsc{ipfv}=2\textsc{sg}.\textsc{pron}\\
\glt `they are what you will use for sewing the sago palm leaves (i.e. the thatch sheets)'
\z

\ea\label{ex:text2-22}
Amaa sikiri bona mee, na veveo raara bona maa nou vaa burihina.\\
\gll a=maa	sikiri	bona	mee	na	veveo	ra=ara	bona	maa	nou	vaa	burihina\\
     \textsc{art}2.\textsc{sg}=\textsc{plm}	midrib	\textsc{ana}.\textsc{sg}	also	\textsc{tm}2	sweep	1\textsc{pl}.\textsc{in}.\textsc{ipfv}=1\textsc{pl}.\textsc{in}.\textsc{pron}	\textsc{obj}.\textsc{art}2.\textsc{sg}	\textsc{plm}	rubbish	\textsc{lnk}	backyard\\
\glt `These midribs we also use for sweeping the rubbish in the backyard.'
\z

\ea\label{ex:text2-23}
O ROM ATOVO\\
\gll o	rom	atovo\\
     \textsc{art}3.\textsc{sg}	stem	sago.palm\\
\glt `THE STEM OF THE SAGO PALM'
\z

\ea\label{ex:text2-24}
O rom atovo na romrom riori, eori repaa kopakopaa bari.\\
\gll o	rom	atovo	na	rom{$\sim$}rom	ri=ori	e-ori	re=paa	kopa{$\sim$}kopaa	bari\\
     \textsc{art}3.\textsc{sg}	stem	sago.palm	\textsc{tm}2	\textsc{red}{$\sim$}cut.crosswise	3\textsc{pl}.\textsc{ipfv}=3\textsc{pl}.\textsc{pron}	\textsc{h}-3\textsc{pl}.\textsc{pron}	\textsc{consec}=\textsc{tm}3	\textsc{red}{$\sim$}cut.lengthwise	4\textsc{pl}.\textsc{pron}\\
\glt `They cut the stem of the sago palm crosswise (into logs), and then they cut them (the logs) lengthwise.'
\z

\ea\label{ex:text2-25}
Murinae eori repaa taraha varaaramun bona komanae bara toon nae,\\
\gll muri-nae	e-ori	re=paa	taraha	va-raa{$\sim$}ramun	bona	koma-nae	bara	toon	nae\\
     after-3\textsc{sg}.\textsc{poss}.\textsc{pron}	\textsc{h}-3\textsc{pl}.\textsc{pron}	\textsc{consec}=\textsc{tm}3	shave	\textsc{caus}-\textsc{red}{$\sim$}smooth	\textsc{obj}.\textsc{art}2.\textsc{sg}	inside-3\textsc{sg}.\textsc{poss}.\textsc{pron}	and2	outside	3\textsc{sg}.\textsc{poss}.\textsc{pron}\\
\glt `After that they shave their insides and outsides to make (them) smooth,'
\z

\ea\label{ex:text2-26}
ahiki ta taba re tahi ribari.\\
\gll ahiki	ta	taba	re	tahi	ri=bari\\
     not.exist	\textsc{nspec}2.\textsc{sg}	thing	\textsc{consec}	prick	3\textsc{pl}.\textsc{objm}=4\textsc{pl}.\textsc{pron}\\
\glt `so that nothing will prick them.'
\z

\ea\label{ex:text2-27}
Beve kavara, eori repaa deedee voosu mana bari vo vaan\\
\gll beve	kavara	e-ori	re=paa	dee{$\sim$}dee	voosu	mana	bari	vo	vaan\\
     when.3\textsc{sg}.\textsc{pron}	finished	\textsc{h}-3\textsc{pl}.\textsc{pron}	\textsc{consec}=\textsc{tm}3	\textsc{red}{$\sim$}carry	go.home	eventually	4\textsc{pl}.\textsc{pron}	\textsc{goal}	village\\
\glt `When it is finished, they eventually carry them home to the village'
\z

\ea\label{ex:text2-28}
mene paku bono baranda na inu.\\
\gll mene	paku	bono	baranda	n=a	inu\\
     for	make	\textsc{obj}.\textsc{art}3.\textsc{sg}	floor	3\textsc{sg}.\textsc{poss}=\textsc{art}2.\textsc{sg}	house\\
\glt `for making the floor of the house.'
\z

\ea\label{ex:text2-29}
VAA'AVE NA INU\\
\gll vaa'ave	n=a	inu\\
     thatch	3\textsc{sg}.\textsc{poss}=\textsc{art}2.\textsc{sg}	house\\
\glt `THATCHING THE HOUSE'
\z

\ea\label{ex:text2-30}
Eori paa huha vakavara vaevuru bona paka atovo vai to paa kosi vuru maori,\\
\gll e-ori	paa	huha	va-kavara	vaevuru	bona	paka	atovo	vai	to	paa	kosi	vuru	ma=ori\\
     \textsc{h}-3\textsc{pl}.\textsc{pron}	\textsc{tm}3	sew	\textsc{advr}-complete	already	\textsc{obj}.\textsc{art}3.\textsc{pl}	leaf	sago.palm	\textsc{dem}6	\textsc{rel}	\textsc{tm}3	cut.off.leaves	earlier	hither=3\textsc{pl}.\textsc{pron}\\
\glt `They have already sewed all (the thatch sheets of) the sago palm leaves that they had cut off earlier,'
\z

\ea\label{ex:text2-31}
matatopo vai tea vaa'ave bona inu.\\
\gll matatopo	vai	tea	vaa'ave	bona	inu\\
     be.prepared	now	\textsc{compl}1	thatch	\textsc{obj}.\textsc{art}2.\textsc{sg}	house\\
\glt `in preparation for thatching the house.'
\z

\ea\label{ex:text2-32}
Evehee beori mei rori tea vaa'ave, toro rahi obai maa roho,\\
\gll evehee	be=ori	mei	rori	tea	vaa'ave	toro	rahi	obai	maa	roho\\
     but	when=3\textsc{pl}.\textsc{pron}	do.not.yet	3\textsc{pl}.\textsc{ipfv}	\textsc{compl}1	thatch	\textsc{tm}7	pull	calamus.cane	hither	before\\
\glt `But before they do the thatching, they must get calamus canes,'\footnote{lit. `But when they have not done yet the thatching, (they) must pull calamus cane hither before.'}
\z

\ea\label{ex:text2-33}
bara toro mosi maa roho ta maa naono vai to pasi sun riori beori tau vaa'ave.\\
\gll bara	toro	mosi	maa	roho	ta	maa	naono	vai	to	pasi	sun	ri=ori	be=ori	tau	vaa'ave\\
     and	\textsc{tm}7	cut	hither	before	\textsc{nspec}2.\textsc{sg}	\textsc{plm}	pole	\textsc{dem}6	\textsc{rel}	\textsc{tm}4	stand	3\textsc{pl}.\textsc{ipfv}=3\textsc{pl}.\textsc{pron}	when=3\textsc{pl}.\textsc{pron}	\textsc{tm}5	thatch\\
\glt `and cut the poles (for the scaffold) on which they will stand when they are going to thatch (the house).'
\z

\ea\label{ex:text2-34}
O dorana o naono pasi vaasun vori,\\
\gll o	dorana	o	naono	pasi	vaa-sun	v=ori\\
     \textsc{art}3.\textsc{sg}	four	\textsc{art}3.\textsc{sg}	pole	\textsc{tm}4	\textsc{caus}-stand	\textsc{pfv}=3\textsc{pl}.\textsc{pron}\\
\glt `They will put up four poles,'
\z

\newpage
\ea\label{ex:text2-35}
eori repaa vaobete vananata bono meho dorana beori kisi batari bari bono kara obai.\\
\gll e-ori	re=paa	va-obete	vananata	bono	meho	dorana	be=ori	kisi	bata=ri	bari	bono	kara	obai\\
     \textsc{h}-3\textsc{pl}.\textsc{pron}	\textsc{consec}=\textsc{tm}3	\textsc{caus}-lie	across	\textsc{obj}.\textsc{art}3.\textsc{sg}	other3.\textsc{sg}	four	when=3\textsc{pl}.\textsc{pron}	tie	\textsc{cont}=3\textsc{pl}.\textsc{ipfv}	4\textsc{pl}.\textsc{pron}	\textsc{obj}.\textsc{art}3.\textsc{sg}	string	calamus\\
\glt `and then they lay another four poles across, while tying them (to the standing poles) with the calamus string.'
\z

\ea\label{ex:text2-36}
Eve a poka bona to pasi sun bata nom ori\\
\gll eve	a	poka	bona	to	pasi	sun	bata	nom	ori\\
     3\textsc{sg}.\textsc{pron}	\textsc{art}2.\textsc{sg}	scaffold	\textsc{ana}.\textsc{sg}	\textsc{rel}	\textsc{tm}4	stand	\textsc{cont}	\textsc{ipfv}	3\textsc{pl}.\textsc{pron}\\
\glt `This is the scaffold on which they will stand,'
\z

\ea\label{ex:text2-37}
beori vaa'ave batari.\\
\gll be=ori	vaa'ave	bata=ri\\
     when=3\textsc{pl}.\textsc{pron}	thatch	\textsc{cont}=3\textsc{pl}.\textsc{ipfv}\\
\glt `when they are doing the thatching.'
\z

\ea\label{ex:text2-38}
Bea poka kavara, taneo vai tea vaa'ave.\\
\gll be=a	poka	kavara	taneo	vai	tea	vaa'ave\\
     when=\textsc{art}2.\textsc{sg}	scaffold	be.finished	start	now	\textsc{compl}1	thatch\\
\glt `When the scaffold is finished, they start thatching.'
\z

\ea\label{ex:text2-39}
A kukan gea dorana aba pasi sunu tea poka,\\
\gll a	kukan	ge=a	dorana	aba	pasi	sun=u	te=a	poka\\
     \textsc{art}2.\textsc{sg}	three	or=\textsc{art}2.\textsc{sg}	four	person	\textsc{tm}4	stand=\textsc{pfv}	\textsc{prep}=\textsc{art}2.\textsc{sg}	scaffold\\
\glt `Three or four people will stand of the scaffold,'
\z

\ea\label{ex:text2-40}
bea meha kukan hee hiava bata ni naori bona vana atovo.\\
\gll be=a	meha	kukan	hee	hiava	bata	ni	nao=ri	bona	vana~atovo\\
     while=\textsc{art}2.\textsc{sg}	other.2\textsc{sg}	three	give	go.up	\textsc{cont}	\textsc{appl}	thither=3\textsc{pl}.\textsc{ipfv}	\textsc{obj}.\textsc{art}3.\textsc{pl}	thatch.sheet\\
\glt `while another three pass the thatch-sheets up (to them).'\footnote{lit. `when another three give the thatch sheets upwards'}
\z

\ea\label{ex:text2-41}
A dorana a aba vai gina to pasi kisi bata ni nom bona vana atovo tea varuvaru na inu.\\
\gll a	dorana	a	aba	vai	gina	to	pasi	kisi	bata	ni	nom	bona	vana~atovo	te=a	varuvaru	n=a	inu\\
     \textsc{art}2.\textsc{sg}	four	\textsc{art}2.\textsc{sg}	person	\textsc{dem}6	on.top	\textsc{rel}	\textsc{tm}4	tie	\textsc{cont}	\textsc{appl}	\textsc{ipfv}	\textsc{obj}.\textsc{art}3.\textsc{pl}	thatch.sheet	\textsc{prep}=\textsc{art}3.\textsc{pl}	rafter	3\textsc{sg}.\textsc{poss}=\textsc{art}2.\textsc{sg}	house\\
\glt `The four people on top (of the scaffold) are the ones that will tie the thatch-sheets to the rafters of the house.'\footnote{lit. `The four people on top (are) who will tie the thatch sheets ...'}
\z

\ea\label{ex:text2-42}
A HUUN NO ATOVO O ANIANI\\
\gll a	huun	n=o	atovo	o	ani{$\sim$}ani\\
     \textsc{art}2.\textsc{sg}	liquid	3\textsc{sg}.\textsc{poss}=\textsc{art}3.\textsc{sg}	sago.palm	\textsc{art}3.\textsc{sg}	\textsc{red}{$\sim$}eat\\
\glt `THE EDIBLE SAGO PALM MASH'
\z

\ea\label{ex:text2-43}
A abana saka tagava pinopino haari bono atovo vai to pasi koro riori mene ani.\\
\gll a	abana	saka	tagava	pinopino	haa=ri	bono	atovo	vai	to	pasi	koro	ri=ori	mene	ani\\
     \textsc{art}2.\textsc{sg}	people	\textsc{neg}1	fell	carelessly	\textsc{neg}2=2\textsc{pl}.\textsc{ipfv}	\textsc{obj}.\textsc{art}3.\textsc{sg}	sago.palm	\textsc{dem}6	\textsc{rel}	\textsc{tm}4	carve.out	3\textsc{pl}.\textsc{ipfv}=3\textsc{pl}.\textsc{pron}	for	eat\\
\glt `The people do not carelessly fell the sago palm that they will carve out for eating (the sago).'
\z

\ea\label{ex:text2-44}
Na tara rori bono atovo vai o vahuusu vua,\\
\gll na	tara	rori	bono	atovo	vai	o	vahuusu	vua\\
     \textsc{tm}2	see	3\textsc{pl}.\textsc{ipfv}	\textsc{obj}.\textsc{art}3.\textsc{sg}	sago.palm	\textsc{dem}6	\textsc{art}3.\textsc{sg}	close.to	bear.nuts\\
\glt `They see a sago palm close to bearing nuts,'
\z

\ea\label{ex:text2-45}
eori repaa tagava bari\\
\gll e-ori	re=paa	tagava	bari\\
     \textsc{h}-3\textsc{pl}.\textsc{pron}	\textsc{consec}=\textsc{tm}3	fell	4\textsc{pl}.\textsc{pron}\\
\glt `and then they fell it,'
\z

\ea\label{ex:text2-46}
ei kou na nata vorien o atovo bona o mataa bara o manoto.\\
\gll ei~kou	na	nata	vo=ri=en	o	atovo	bona	o	mataa	bara	o	manoto\\
     because	\textsc{tm}2	know	\textsc{goal}=3\textsc{pl}.\textsc{ipfv}=\textsc{dem}3	\textsc{art}3.\textsc{sg}	sago.palm	\textsc{ana}.\textsc{sg}	\textsc{art}3.\textsc{sg}	good	and2	\textsc{art}3.\textsc{sg}	ripe\\
\glt `because they know that this sago palm is good and ripe.'
\z

\ea\label{ex:text2-47}
Beori tagava bari, eori repaa kopaa pete bari.\\
\gll be=ori	tagava	bari	e-ori	re=paa	kopaa	pete	bari\\
     when=3\textsc{pl}.\textsc{pron}	fell	4\textsc{pl}.\textsc{pron}	\textsc{h}-3\textsc{pl}.\textsc{pron}	\textsc{consec}=\textsc{tm}3	cut.lengthwise	subsequently	4\textsc{pl}.\textsc{pron}\\
\glt `When they have felled it, they cut it lengthwise.'
\z

\ea\label{ex:text2-48}
Murinae eori repaa koro bari bara tasu vamuramuraka bata bari.\\
\gll muri-nae	e-ori	re=paa	koro	bari	bara	tasu	va-mura{$\sim$}muraka	bata	bari\\
     after-3\textsc{sg}.\textsc{poss}.\textsc{pron}	\textsc{h}-3\textsc{pl}.\textsc{pron}	\textsc{consec}=\textsc{tm}3	carve.out	4\textsc{pl}.\textsc{pron}	and2	beat	\textsc{caus}-\textsc{red}{$\sim$}soft	\textsc{cont}	4\textsc{pl}.\textsc{pron}\\
\glt `After that they cut it out and mash it (i.e. what they have carved out).'
\z

\ea\label{ex:text2-49}
Eori repaa garuhu bari ruene.\\
\gll e-ori	re=paa	garuhu	bari	ruene\\
     \textsc{h}-3\textsc{pl}.\textsc{pron}	\textsc{consec}=\textsc{tm}3	wash	4\textsc{pl}.\textsc{pron}	fresh.water\\
\glt `Then they wash it (the mash) in fresh water.'
\z

\ea\label{ex:text2-50}
Beori garuhu bari ore paa posi batari bari vo tea maa nahu vai amaa beera ge pagete.\\
\gll be=ori	garuhu	bari	o=re	paa	posi	bata=ri	bari	vo	te=a	maa	nahu	vai	a=maa	beera	ge	pagete\\
     when=3\textsc{pl}.\textsc{pron}	wash	4\textsc{pl}.\textsc{pron}	3.\textsc{pron}=\textsc{consec}	\textsc{tm}3	squeeze	\textsc{cont}=3\textsc{pl}.\textsc{ipfv}	4\textsc{pl}.\textsc{pron}	\textsc{goal}	\textsc{prep}=\textsc{art}2.\textsc{sg}	\textsc{plm}	pot	\textsc{dem}6	\textsc{art}2.\textsc{sg}=\textsc{plm}	large	or	bucket\\
\glt `When they have washed it, they squeeze it (the water out of the mash) into large pots or buckets.'
\z

\newpage
\ea\label{ex:text2-51}
O hum vai to ani riori pasi taviana vaavaha komana bata nom.\\
\gll o	hum	vai	to	ani	ri=ori	pasi	ta-viana	vaavaha	komana	bata	nom\\
     \textsc{art}3.\textsc{sg}	part	\textsc{dem}6	\textsc{rel}	eat	3\textsc{pl}.\textsc{ipfv}=3\textsc{pl}.\textsc{pron}	\textsc{tm}4	\textsc{acaus}-select	separate	self	\textsc{cont}	\textsc{ipfv}\\
\glt `The part that they eat gets separated.'
\z

\ea\label{ex:text2-52}
Beve raka, ore paa kakaavoo.\\
\gll beve	raka	o=re	paa	kakaavoo\\
     when.3\textsc{sg}.\textsc{pron}	dry	3.\textsc{pron}=\textsc{consec}	\textsc{tm}3	white\\
\glt `When it is dry, it becomes white.'
\z

\ea\label{ex:text2-53}
MOROKO VAGEGEANA\\
\gll moroko	vagegeana\\
     talk	last\\
\glt `CONCLUSION'
\z

\ea\label{ex:text2-54}
O atovo o mataa, a haihai teve a beera sana en teara to vigu inu ni rara o atovo.\\
\gll o	atovo	o	mataa	a	haihai	teve	a	beera	sana	en	te=ara	to	vigu	inu	ni	rara	o	atovo\\
     \textsc{art}3.\textsc{sg}	sago.palm	\textsc{art}3.\textsc{sg}	good	\textsc{art}2.\textsc{sg}	help	\textsc{prep}.3\textsc{sg}.\textsc{pron}	\textsc{art}2.\textsc{sg}	big	very	\textsc{dem}3	\textsc{prep}=1\textsc{pl}.\textsc{in}.\textsc{pron}	\textsc{rel}	build	house	\textsc{appl}	1\textsc{pl}.\textsc{in}.\textsc{ipfv}	\textsc{art}3.\textsc{sg}	sago.palm\\
\glt `The sago palm is good, its help is very big for us who use the sago palm for building houses.'
\z
